\begin{frame}{Why familiar algebraic laws survive}
  \begin{columns}[T,onlytextwidth]
    \begin{column}{0.54\textwidth}
      \[
      A+B=B+A
      \]
      \[
      (A+B)+C=A+(B+C)
      \]
      \[
      \lambda(A+B)=\lambda A+\lambda B
      \]
    \end{column}
    \begin{column}{0.42\textwidth}
      \begin{exampleblock}{Reason}
        Each statement is true entry by entry because ordinary numbers satisfy the same laws.
      \end{exampleblock}
    \end{column}
  \end{columns}
\end{frame}

\begin{frame}{Before matrix multiplication}
  \begin{alertblock}{Entrywise model}
    This model belongs to addition and scalar multiplication.
  \end{alertblock}

  \vspace{0.5cm}
  \begin{exampleblock}{Row--column model}
    Matrix multiplication makes a row interact with a column.
  \end{exampleblock}

  \vspace{0.55cm}
  \concept{The next operation changes the logic, not only the formula.}
\end{frame}

\begin{frame}{Checkpoint: entrywise operations}
  \begin{columns}[T,onlytextwidth]
    \begin{column}{0.62\textwidth}
      \Large
      \begin{itemize}
        \item same size is required;
        \item shape is preserved;
        \item every entry keeps its address;
        \item familiar arithmetic laws pass through.
      \end{itemize}
    \end{column}
    \begin{column}{0.34\textwidth}
      \takeaway{Position is the organizing principle.}
    \end{column}
  \end{columns}
\end{frame}
